
%\chapter{Introdução}

%\textbf{Organização de informação não se dá da mesma forma em qualquer domínio.}
O mundo passou e vem passando por grandes mudanças tecnologicas nos ultimos anos, o que possibitou a presença da tecnologia em diversos setores econômicos \cite{del2009avancco}. Juntamente com a presença tecnologica as pessoas puderam se tornar mais ativos virtualmente, uma vez que a tecnologia faz parte do dia a das pessoas e em conjunto com essa presença digital se tornou notável o consumo de conteúdos de outros páises




\section{Justificativa} %Inseri para organizar o texto

Hoje em dia existem diversos modelos neurais de tradução de código aberto, porém cada modelo possui dados de trinamento diferentes, e por mais que todos os modelos presentados neste trabalho sejam baseados na arquietura transformer, cada modelo terá caracteristicas arquiteturais que irão diferenciar um modelo do oytro, impactando o resultado final da tarefa de tradução.Este trabalho visa comaparara cada modelo a fim de mostrar onde e porque cada modelo se sobressaí em relação ao outro quando se trata de uma tarefa de tradução. 

\section{Problema}

A seguinte questão constitue o problema desta pesquisa: 
Qual é o desempenho de diferentes modelos de tradução entre línguas e automática, baseados na arquitetura Transformer para a tarefa de tradução do inglês para o português, considerando métricas como consumo de memória, uso de CPU, pontuação BLEU, BERTScore e chrF, em diferentes cenários de teste?


\section{Objetivos}

%\textbf{Objetivo geral.}
Para responder ao problema proposto, o objetivo geral deste trabalho é investigar e comparar o desempenho de modelos de tradução entre línguas e automática baseados na arquitetura Transformer na tarefa de tradução do inglês para o português, considerando métricas de qualidade (BLEU, BERTScore e chrF) e eficiência computacional (uso de memória e CPU) em diferentes cenários experimentais.

Também, objetivam-se mais especificamente:

\begin{enumerate}
	\item Revisão bibliográfica sobre o funcionamento de modelos de tradução automática neural baseados na arquitetura Transformer;

	\item Selecionar e configurar os modelos de tradução automática neural de código aberto;

	\item Selecionar base de dados pública para comparação entre os modelos de tradução automática neural;
	
    \item Avaliar qualidade de tradução utilizando metricas como BLEU, BERTScore e chrF;
    
    \item Medir eficiência de recursos computacionais quando se utiliza modelos de tradução;
    
    \item Avaliar os modelos de tradução diferentes cenários de teste;
    
    \item Analisar correlação entre métricas de qualidade de tradução.

\end{enumerate}



